\chapter{Compliance and Market Abuse}\label{ch:fm:compliance-and-market-abuse}

In August 2022 the UK's Financial Conduct Authority fined Citigroup Global Markets \pounds12.6 million because, when the Market Abuse Regulation
extended the duty to monitor orders as well as trades in 2016, the firm did not properly implement it: it took eighteen months to identify and
assess the \omterm{def:fm:compliance-and-market-abuse:abuse}{market abuse} risks its business was exposed to, and in the meantime it could not effectively monitor its trading for certain types of
\omterm{def:fm:compliance-and-market-abuse:abuse}{insider dealing} and market manipulation. Nobody was accused of abusing the market. The failure was significant gaps in the firm's surveillance: for some kinds of
abuse no alert would have been generated, so nobody could have read one. Compliance is a function that must see the trading before it can
judge it.

\section{The rulebook a trader lives under}

A trader's conduct is bounded by two prohibitions and the duties that support them. The prohibitions are on using information others do not have
and on distorting prices; Book 9 described the manipulative patterns (spoofing, layering, marking the close, momentum ignition, benchmark
manipulation) and Book 3 wash trading. This chapter is about the function that keeps a firm on the right side of them.

\begin{definition}[Market abuse, inside information, insider dealing]\label{def:fm:compliance-and-market-abuse:abuse}
\emph{Market abuse}\index{market abuse} is insider dealing, the unlawful disclosure of inside information, and market manipulation, including
attempts. \emph{Inside information}\index{inside information} is information of a precise nature, not made public, relating directly or indirectly
to an issuer or a financial instrument, which would be likely to have a significant effect on prices if it were made public. \emph{Insider
dealing}\index{insider dealing} is using inside information to buy or sell, for oneself or another, the instruments to which it relates, including
by cancelling or amending an order placed before the information was received.
\end{definition}

\begin{dated}{2026-09}{The EU Market Abuse Regulation}\label{dat:fm:compliance-and-market-abuse:mar}
Regulation (EU) No 596/2014 defines \omterm{def:fm:compliance-and-market-abuse:abuse}{inside information} (article 7) and \omterm{def:fm:compliance-and-market-abuse:abuse}{insider dealing} (article 8); prohibits \omterm{def:fm:compliance-and-market-abuse:abuse}{insider dealing}, recommending or
inducing it and unlawfully disclosing \omterm{def:fm:compliance-and-market-abuse:abuse}{inside information} (article 14), and market manipulation (article 15); governs market soundings, whose
records are kept for at least five years (article 11); requires any person professionally arranging or executing transactions to maintain
effective arrangements, systems and procedures to detect and report suspicious orders and transactions, and to notify the competent authority
without delay on reasonable suspicion (article 16(2)); and requires issuers and those acting for them to keep \omterm{def:fm:compliance-and-market-abuse:barrier}{insider lists} (article 18).
\end{dated}

\begin{dated}{2026-09}{US rules on insider trading and personal dealing}\label{dat:fm:compliance-and-market-abuse:us}
Rule 10b-5 (17 CFR 240.10b-5) makes it unlawful to employ any device, scheme or artifice to defraud, or to engage in any act that operates as a
fraud, in connection with the purchase or sale of any security; insider trading cases are brought under it. Section 15(g) of the Exchange Act (15
U.S.C. 78o(g)) requires every registered broker-dealer to establish, maintain and enforce written policies and procedures reasonably designed to
prevent the misuse of material non-public information. Rule 204A-1 (17 CFR 275.204A-1) requires a registered investment adviser's code of ethics
to make its access persons report their holdings and, within 30 days after each quarter, their personal transactions, and to obtain approval
before buying in an initial public offering or a limited offering.
\end{dated}

\section{Surveillance as a function}

Book 15 built the trade surveillance system: detectors that score orders and trades, surveillance alerts, triage, and the communications archive.
The compliance function owns what happens next: which detectors run on which businesses, where their thresholds sit, who reviews the alerts, and
what is reported.

\begin{definition}[Suspicious transaction and order report]\label{def:fm:compliance-and-market-abuse:stor}
A \emph{suspicious transaction and order report}\index{suspicious transaction and order report} (STOR) is the notification a firm that arranges or
executes transactions must send its regulator, without delay, when it has a reasonable suspicion that an order or transaction could be \omterm{def:fm:compliance-and-market-abuse:abuse}{insider
dealing} or market manipulation, attempted or completed.
\end{definition}

The hook's failure was coverage: detectors that were not run cannot alert. The subtler failure is capacity. Every detector has a threshold; below
it nothing is reviewed, above it every alert waits for an analyst. Set it low and the queue floods, and analysts reviewing in the order alerts arrive
see mostly noise; set it high and true cases never alert.

\begin{proposition}[Filling the queue]\label{prop:fm:compliance-and-market-abuse:queue}
Let a detector score $N$ candidates, and let an alert rate $q$ send the $qN$ highest-scored for review, of which $\mathrm{TP}(q)$ are true cases,
non-decreasing in $q$, with a precision $\mathrm{TP}(q)/(qN)$ non-increasing in $q$. If analysts review at most $C$ alerts, in arrival order, the
expected number of true cases found is $\mathrm{TP}(q)$ for $q\le C/N$ and $C\,\mathrm{TP}(q)/(qN)$ above it, and is largest at $q^*=C/N$: the
threshold whose alerts just fill the analysts' capacity.
\end{proposition}

\begin{proof}
Below $C/N$ every alert is reviewed and the true cases found are $\mathrm{TP}(q)$, which does not fall as $q$ rises. Above it a share $C/(qN)$ of
alerts is reviewed, in an order unrelated to the score, so the expected true cases found are $C$ times the precision, which does not rise with
$q$. Both pieces equal $\mathrm{TP}(C/N)$ at $q=C/N$.
\end{proof}

\omcode{code/firm/compliance/firm_compliance.py}{76}{93}{The triage queue: alerts above the threshold for a chosen alert rate, reviewed up to
capacity in arrival or score order, and the alert rate that finds the most true cases.}\label{lst:fm:compliance-and-market-abuse:triage}

\section{Tutorial: the alert queue and the pre-clearance desk}\label{tut:fm:compliance-and-market-abuse:tut}

\begin{tutorial}
\textbf{Goal.} Run a year of personal-trade \omterm{def:fm:compliance-and-market-abuse:pad}{pre-clearance}, then send a surveillance alert stream through a triage queue with fixed analyst hours
and find the threshold that catches the most true cases. \textbf{End state:} the \omterm{def:fm:compliance-and-market-abuse:pad}{pre-clearance} decisions by reason and the queue chart
(\cref{fig:fm:compliance-and-market-abuse:queue}).
\begin{tutsteps}
\item \textbf{The alerts.} Book 9's synthetic spoofing data, \texttt{firm.surveil.spoofing\_days}: 11\,600 account-days, 100 of them planted
abuse, scored by its best detector (large-order fill gap times cancellations after the other side fills).
\item \textbf{The analysts.} Each alert takes half an hour; the team has 50 hours for this detector (100 alerts), or 150 (300 alerts).
\item \textbf{The queue.} \texttt{fm\_compliance.queue(hours)} sweeps the alert rate from 0.1\% to 20\% with \texttt{firm.compliance.triage}.
\item \textbf{The requests.} \texttt{fm\_compliance.year()} sends 1\,215 personal-trade requests from 150 employees through
\texttt{firm.compliance.preclear} against lists that change through the year.
\end{tutsteps}
\end{tutorial}

\begin{omfigure}
\begin{tikzpicture}
\begin{axis}[width=10.5cm,height=5.4cm,xmode=log,log basis x=10,xtick={0.1,0.2,0.5,1,2,5,10,20},xticklabels={0.1,0.2,0.5,1,2,5,10,20},xlabel={\small alert rate (\% of account-days, log scale)},
  ylabel={\small true cases found (of 100)},xmin=0.1,xmax=20,ymin=0,ymax=100,
  tick label style={font=\footnotesize},grid=major,grid style={gray!20},table/col sep=comma,
  legend cell align=left,legend style={font=\scriptsize,at={(0.5,-0.3)},anchor=north,/tikz/every even column/.append style={column sep=8pt}},legend columns=2]
\addplot[omDef,thick,mark=*,mark size=1pt] table[x=rate,y=h50]{figdata/desk/16-compliance-and-market-abuse/queue.csv};
\addplot[omThm,thick,dashed,mark=square*,mark size=1pt] table[x=rate,y=h150]{figdata/desk/16-compliance-and-market-abuse/queue.csv};
\draw[gray,dotted] (axis cs:5,0) -- (axis cs:5,100);
\node[font=\scriptsize,anchor=south west] at (axis cs:5.3,82) {default 5\%};
\legend{50 analyst hours,150 analyst hours}
\end{axis}
\end{tikzpicture}

\omcaption{True cases found by an alert queue reviewed in arrival order, against the alert rate, for two analyst budgets. Each curve peaks where the
alerts just fill the capacity (\cref{prop:fm:compliance-and-market-abuse:queue}); at the default rate of 5\% the queue floods and most true alerts
are never read. Data: \texttt{fm\_compliance.queue}.}\label{fig:fm:compliance-and-market-abuse:queue}
\end{omfigure}

With 50 analyst hours the best alert rate is 0.8\%, just under the capacity rate of $100/11\,600=0.86\%$, and the queue finds 65 of the 100
planted cases. At the default rate of 5\% the detector alerts on 580 account-days, 89 of them true, but analysts reading the first 100 in arrival
order find 15 on average: 85\% of the abuse is missed, not by the detector, which flagged most of it, but by the queue
(\cref{fig:fm:compliance-and-market-abuse:funnel}). Tripling the hours moves the best rate to 2.5\% and the true cases found to 81; at the default
rate it finds 46. Reviewing in score order instead of arrival order recovers what the threshold loses: 66 found at the default rate with 50 hours,
81 with 150. A naive detector, the order-to-trade ratio, finds none of the 100 at any threshold, because the market makers of Book 9 cancel more
than any spoofer.

\begin{omfigure}
\begin{tikzpicture}[font=\footnotesize,>=stealth,box/.style={draw,rounded corners,align=center,minimum height=0.8cm}]
\node[box,fill=gray!15,text width=10cm] (a) at (0,3.0) {11\,600 account-days scored, 100 planted cases};
\node[box,fill=omDef!15,text width=7.6cm] (b) at (0,1.8) {580 alerts at the default 5\% rate, 89 true};
\node[box,fill=omProp!15,text width=5cm] (c) at (0,0.6) {100 reviewed in 50 hours};
\node[box,fill=omThm!15,text width=2.8cm] (d) at (0,-0.6) {15 true found};
\draw[->] (a) -- (b); \draw[->] (b) -- (c); \draw[->] (c) -- (d);
\node[font=\scriptsize,align=left,anchor=west] at (3.9,1.8) {detector: 89\% recall};
\node[font=\scriptsize,align=left,anchor=west] at (2.7,0.6) {queue: 17\% of alerts read};
\end{tikzpicture}

\omcaption{The surveillance funnel at the default threshold: the detector finds most of the abuse and the queue loses it. Data:
\texttt{fm\_compliance.queue}.}\label{fig:fm:compliance-and-market-abuse:funnel}
\end{omfigure}

\section{Personal dealing}

Employees who know what the firm is about to trade, or what a client is about to announce, can profit from it in their own accounts. The rules are
older than algorithmic trading and apply to everyone with access.

\begin{definition}[Personal account dealing, pre-clearance]\label{def:fm:compliance-and-market-abuse:pad}
\emph{Personal account dealing}\index{personal account dealing} is trading by employees for their own accounts or accounts they control or
benefit from. \emph{Pre-clearance}\index{pre-clearance} is the requirement that an employee obtain the compliance function's approval before a
personal trade, which is refused if the instrument is restricted, if the firm is trading or about to trade it, if a minimum holding period has not
passed, or if the employee holds \omterm{def:fm:compliance-and-market-abuse:abuse}{inside information} about it.
\end{definition}

\begin{definition}[Restricted list, watch list]\label{def:fm:compliance-and-market-abuse:lists}
A \emph{restricted list}\index{restricted list} names the instruments in which the firm and its employees may not trade, or may trade only with
approval, usually because the firm holds \omterm{def:fm:compliance-and-market-abuse:abuse}{inside information} about the issuer and says so publicly or to its staff. A \emph{watch
list}\index{watch list} names, confidentially and for compliance only, the instruments about which some part of the firm holds \omterm{def:fm:compliance-and-market-abuse:abuse}{inside information},
so that trading in them anywhere in the firm can be monitored without revealing the information.
\end{definition}

\omcode{code/firm/compliance/firm_compliance.py}{41}{57}{Pre-clearance of a personal trade: the first rule that applies gives the reason for a
refusal.}\label{lst:fm:compliance-and-market-abuse:preclear}

\begin{omfigure}
\begin{tikzpicture}
\begin{axis}[width=10.5cm,height=4.8cm,xbar,bar width=9pt,xmin=0,xmax=1100,xlabel={\small requests},ytick={0,1,2,3,4,5},
  yticklabels={approved,firm traded recently,restricted list,insider,holding period,watch list},y dir=reverse,ymin=-0.6,ymax=5.6,
  tick label style={font=\footnotesize},grid=major,grid style={gray!20},table/col sep=comma,nodes near coords,
  nodes near coords style={font=\scriptsize}]
\addplot[fill=omDef!55,draw=omDef] table[y=pos,x=count]{figdata/desk/16-compliance-and-market-abuse/reasons.csv};
\end{axis}
\end{tikzpicture}

\omcaption{A year of personal-trade requests by decision: 1\,215 requests, 964 approved; four in five refusals are for names the firm traded in
the two days before. Data: \texttt{fm\_compliance.reasons}.}\label{fig:fm:compliance-and-market-abuse:reasons}
\end{omfigure}

Of the year's 1\,215 requests, 964 (79.3\%) were approved. Of the 251 refusals, 202 (80.5\%) were for names the firm had traded in the previous two
days, the front-running rule; 23 were for restricted names, 5 for names on the \omterm{def:fm:compliance-and-market-abuse:lists}{watch list}, 6 for sales within the 30-day holding period, and 15
were requests by wall-crossed employees in the names of their projects, every one of the planted requests refused. The watch-list refusals show
the list's cost: an employee refused for a reason compliance cannot give learns that something is happening in the name. A firm can reduce that
signal by giving no reasons for any refusal, or by monitoring watch-list names instead of refusing trades in them.

\section{Information barriers, insider lists and wall-crossing}

A firm that advises issuers, lends to them or underwrites their securities holds \omterm{def:fm:compliance-and-market-abuse:abuse}{inside information} in one part and trades in another. It keeps
them apart with barriers, and it records who crosses them.

\begin{definition}[Information barrier, wall-crossing, insider list]\label{def:fm:compliance-and-market-abuse:barrier}
An \emph{information barrier}\index{information barrier} is the set of organisational, physical and electronic controls that keep \omterm{def:fm:compliance-and-market-abuse:abuse}{inside
information} held in one part of a firm from people in another who trade or advise on the instruments concerned. A
\emph{wall-crossing}\index{wall-crossing} is the controlled passing of \omterm{def:fm:compliance-and-market-abuse:abuse}{inside information} to a person on the other side of a barrier, or to an
outside investor sounded about a transaction, with their prior consent, recorded, and binding them not to trade until the information is public or
no longer inside. An \emph{insider list}\index{insider list} records every person with access to \omterm{def:fm:compliance-and-market-abuse:abuse}{inside information} about an issuer, when they
obtained it, and why.
\end{definition}

\begin{omfigure}
\begin{tikzpicture}[font=\footnotesize,>=stealth,box/.style={draw,rounded corners,align=center,minimum height=0.8cm,text width=3cm}]
\node[box,fill=omDef!15] (priv) at (-3.8,0) {private side\\{\scriptsize advisory, lending, capital markets}};
\node[box,fill=omProp!15] (pub) at (3.8,0) {public side\\{\scriptsize trading, research, sales}};
\draw[very thick,gray] (0,-1.4) -- (0,1.6) node[above,font=\scriptsize,text=black] {information barrier};
\node[box,fill=omThm!15,text width=3.4cm] (comp) at (0,-2.6) {compliance\\{\scriptsize watch list, insider list, restricted list}};
\draw[->,dashed] (priv) -- (pub); \node[font=\scriptsize] at (-1.1,0.28) {wall-crossing};
\draw[->] (priv.south) |- (comp.west); \draw[->] (comp.east) -| (pub.south);
\node[font=\scriptsize,align=left,anchor=west] at (3.95,-1.3) {pre-clearance,\\surveillance};
\node[font=\scriptsize,anchor=east] at (-3.95,-1.3) {deals logged};
\end{tikzpicture}

\omcaption{\omterm{def:fm:compliance-and-market-abuse:barrier}{Information barriers}: the private side holds \omterm{def:fm:compliance-and-market-abuse:abuse}{inside information}, the public side trades; compliance sits across both, keeps the lists,
records \omterm{def:fm:compliance-and-market-abuse:barrier}{wall-crossings} and pre-clears and surveils the public side. Schematic.}\label{fig:fm:compliance-and-market-abuse:barrier}
\end{omfigure}

A crossed person is on the \omterm{def:fm:compliance-and-market-abuse:barrier}{insider list} for the project's issuer from the moment they are crossed until the information is made public or loses
its significance (the project is cleansed); \omterm{def:fm:compliance-and-market-abuse:pad}{pre-clearance} refuses their trades in the name throughout, and surveillance watches the name across the
firm. The records serve two purposes: they make the controls work, and they are what the firm shows its regulator when a trade in the name is
questioned. A market sounding under the regulation's rules is the same act performed on an outside investor, with the same consent and records.

\begin{method}[Running the compliance function]\label{met:fm:compliance-and-market-abuse:run}
\begin{enumerate}
\item Map the firm's businesses to the abuse each could commit or facilitate, and each to a detector; check the coverage every time a business,
product or venue is added.
\item Set each detector's threshold so that its alerts fill the analysts' hours, and review in score order; measure the true-alert rate and adjust.
\item Keep the lists point in time; pre-clear every personal trade; record every \omterm{def:fm:compliance-and-market-abuse:barrier}{wall-crossing} and cleansing.
\item Report suspicious orders and transactions without delay, and record the reasons for every decision not to report.
\item Test the whole chain once a year with planted cases, as the tutorial does.
\end{enumerate}
\end{method}

\section{Build: the compliance engine}\label{bld:fm:compliance-and-market-abuse:compliance}

\begin{build}
\bfield{Purpose} Restricted and \omterm{def:fm:compliance-and-market-abuse:lists}{watch lists} through time, \omterm{def:fm:compliance-and-market-abuse:pad}{pre-clearance} of personal trades, an insider-list register with \omterm{def:fm:compliance-and-market-abuse:barrier}{wall-crossings}, and the
capacity of an alert queue on firm.surveil's detectors.
\bfield{Interface} \texttt{firm.compliance}: \texttt{Lists}, \texttt{Rules}, \texttt{preclear}, \texttt{Register} (\texttt{cross},
\texttt{cleanse}, \texttt{insiders}), \texttt{triage}, \texttt{best\_rate}.
\bfield{Rules} The first applicable rule decides a refusal and is recorded; lists and registers are dated; a cleansed project's insiders leave the
list on the cleansing day; triage never reviews more alerts than capacity.
\bfield{Acceptance tests} \texttt{code/firm/compliance/tests/}: each refusal reason; the register through a project; the queue on a detector with
known precision.
\bfield{Stretch} Alerts from several detectors sharing one team; review times that depend on the alert; the regulator's reporting deadline as a
constraint on the queue.
\end{build}

\begin{omsources}
\item Financial Conduct Authority, press release of 19 August 2022 on Citigroup Global Markets Limited.
\item Regulation (EU) No 596/2014, articles 7, 8, 11, 14--16 and 18.
\item 17 CFR 240.10b-5; 15 U.S.C. 78o(g); 17 CFR 275.204A-1.
\end{omsources}

\section{Exercises}

\begin{exercise}[$\star$]\label{exo:fm:compliance-and-market-abuse:1}
Which of the chapter's refusal rules would stop a trader who wants to buy a stock the firm bought yesterday, and why does the rule exist?
\end{exercise}

\begin{exercise}[$\star$]\label{exo:fm:compliance-and-market-abuse:2}
A detector scores 11\,600 account-days and analysts can review 100 alerts. What alert rate fills the queue?
\end{exercise}

\begin{exercise}[$\star$]\label{exo:fm:compliance-and-market-abuse:3}
What is the difference between a \omterm{def:fm:compliance-and-market-abuse:lists}{restricted list} and a \omterm{def:fm:compliance-and-market-abuse:lists}{watch list}, and why is the \omterm{def:fm:compliance-and-market-abuse:lists}{watch list} confidential?
\end{exercise}

\begin{exercise}[$\star\star$]\label{exo:fm:compliance-and-market-abuse:4}
At the default 5\% rate the detector flags 89 of the 100 cases. Why does the queue find only 15?
\end{exercise}

\begin{exercise}[$\star\star$]\label{exo:fm:compliance-and-market-abuse:5}
Why does the order-to-trade ratio find none of the planted spoofing cases?
\end{exercise}

\begin{exercise}[$\star\star$]\label{exo:fm:compliance-and-market-abuse:6}
A research analyst is wall-crossed on a takeover on day 40 and the project is cleansed on day 75. On which days is a personal trade in the target
refused, and why is surveillance of the whole firm's trading in the name still needed?
\end{exercise}

\begin{exercise}[$\star\star\star$]\label{exo:fm:compliance-and-market-abuse:7}
\emph{Coding.} Run \texttt{fm\_compliance.queue} with 150 hours, in arrival order and in score order. What do the best rates and the default
rate find?
\end{exercise}

\begin{exercise}[$\star\star\star$]\label{exo:fm:compliance-and-market-abuse:8}
\emph{Find the flaw.} ``We lowered every threshold to catch more abuse; our surveillance is stronger than last year.''
\end{exercise}

\section{Problem: The Alert Queue}

\begin{problem}[{Weekend problem --- the alert queue}]\label{pb:fm:compliance-and-market-abuse:1}
A new head of compliance inherits a surveillance team that clears its queue every quarter by closing old alerts unread, and a \omterm{def:fm:compliance-and-market-abuse:pad}{pre-clearance} process
run by email.

\medskip\noindent\textbf{Part I --- The rules.}
\begin{enumerate}
\item Define \omterm{def:fm:compliance-and-market-abuse:abuse}{market abuse}, \omterm{def:fm:compliance-and-market-abuse:abuse}{inside information} and \omterm{def:fm:compliance-and-market-abuse:abuse}{insider dealing}.
\item What does the EU regulation require of firms that arrange or execute transactions?
\item State the US rules on insider trading, broker-dealer policies and advisers' codes of ethics.
\item What failure did the FCA find at Citigroup Global Markets in 2022?
\end{enumerate}

\medskip\noindent\textbf{Part II --- The queue.}
\begin{enumerate}[resume]
\item Define a STOR.
\item State and prove \cref{prop:fm:compliance-and-market-abuse:queue}.
\item Give the best rates and true cases found for 50 and 150 analyst hours.
\item Give what the default rate finds and misses with 50 hours.
\item What does reviewing in score order change?
\item Why is a detector's recall not the function's recall?
\end{enumerate}

\medskip\noindent\textbf{Part III --- Personal dealing and barriers.}
\begin{enumerate}[resume]
\item Define \omterm{def:fm:compliance-and-market-abuse:pad}{personal account dealing}, \omterm{def:fm:compliance-and-market-abuse:pad}{pre-clearance}, the \omterm{def:fm:compliance-and-market-abuse:lists}{restricted list} and the \omterm{def:fm:compliance-and-market-abuse:lists}{watch list}.
\item Give the year's \omterm{def:fm:compliance-and-market-abuse:pad}{pre-clearance} decisions by reason.
\item Define an \omterm{def:fm:compliance-and-market-abuse:barrier}{information barrier}, a \omterm{def:fm:compliance-and-market-abuse:barrier}{wall-crossing} and an \omterm{def:fm:compliance-and-market-abuse:barrier}{insider list}.
\item How long does a crossed person stay on the \omterm{def:fm:compliance-and-market-abuse:barrier}{insider list}, and what follows for their trades?
\end{enumerate}

\medskip\noindent\textbf{Part IV --- The plan.}
\begin{enumerate}[resume]
\item What is wrong with closing old alerts unread?
\item How would you set each detector's threshold?
\item How would you check coverage when the firm starts trading a new product?
\item How would you test the whole chain?
\item State the \emph{named result}: the alert threshold that maximises true cases found for a fixed number of analyst hours, and the share
missed at the default threshold.
\item In two sentences, write the plan.
\end{enumerate}
\end{problem}

\section{Interview questions}

\begin{interviewq}[$\star$ \iqroles{trader}]\label{iq:fm:compliance-and-market-abuse:1}
What is \omterm{def:fm:compliance-and-market-abuse:abuse}{inside information}, and what must you do if you receive some by mistake?
\end{interviewq}

\begin{interviewq}[$\star$ \iqroles{trader, researcher}]\label{iq:fm:compliance-and-market-abuse:2}
Why does your firm make you pre-clear personal trades?
\end{interviewq}

\begin{interviewq}[$\star\star$ \iqroles{developer}]\label{iq:fm:compliance-and-market-abuse:3}
Design the data model for a \omterm{def:fm:compliance-and-market-abuse:lists}{restricted list} that can answer ``was this name restricted at 10:32 on 3 March?''
\end{interviewq}

\begin{interviewq}[$\star\star$ \iqroles{researcher}]\label{iq:fm:compliance-and-market-abuse:4}
A surveillance detector has 90\% recall at a 5\% alert rate. Is it good?
\end{interviewq}

\begin{interviewq}[$\star\star$ \iqroles{risk}]\label{iq:fm:compliance-and-market-abuse:5}
What is a \omterm{def:fm:compliance-and-market-abuse:barrier}{wall-crossing}, and what records should it leave?
\end{interviewq}

\begin{interviewq}[$\star\star\star$ \iqroles{researcher, developer}]\label{iq:fm:compliance-and-market-abuse:6}
Several detectors share one team of analysts. How would you allocate their hours?
\end{interviewq}
